\ifdefined\gkver\else\def\gkver{student}\fi
\documentclass[\gkver]{gaokaozhenti}
\begin{document}

\chapter{2004年全国理综Ⅲ}
%% paperType: 理综物理部分

\begin{enumerate}
\item
%% number: 15
%% paperName: 2004年全国理综Ⅲ第 15 题 6 分
%% typeId: 1
%% type: 单选题
%% chapter: 原子和原子核
%% point: 质能方程
%% method: 无
%% score: 6
%% degree: 850
%% duplicateId: 0
%% body:
以 $m_{D}$、$m_{P}$、$m_{n}$ 分别表示氘核、质子、中子的质量，则\xzanswer{D}
\fourchoices[answer=D]
{$m_{D}=m_{p}+m_{n}$}
{$m_{D}=m_{p}+2m_{n}$}
{$m_{D}>m_{p}+m_{n}$}
{$m_{D}<m_{P}+m_{n}$}
%% answer: D
%% memo:
\memoanswer{
质子与电子聚变生成氘核过程中有质量亏损，释放能量，所以 $m_{D}<m_{p}+m_{n}$，故 D 正确。
}

\item
%% number: 16
%% paperName: 2004年全国理综Ⅲ第 16 题 6 分
%% typeId: 1
%% type: 单选题
%% chapter: 机械波
%% point: 波的多解性
%% method: 无
%% score: 6
%% degree: 750
%% duplicateId: 0
%% body:
一简谐横波在图中 $x$ 轴上传播，实线和虚线分别是 $t_{1}$ 和 $t_{2}$ 时刻的波形图，已知 $t_{2}-t_{1}=1.0\Us$。由图判断下列哪一个波速是不可能的\xzanswer{D}
\onepicture[w=9.0cm]{figs/16.png}
\fourchoices[answer=D]
{$1\Ums$}
{$3\Ums$}
{$5\Ums$}
{$10\Ums$}
%% answer: D
%% memo:
\memoanswer{
由图可知波长 $\lambda=4\Um$。若波向右传播，$v=\frac{(n+\frac{1}{4})\lambda}{t_{2}-t_{1}}=(4n+1)\Ums$（$n=0,1,2,\ldots$），$v$ 的可能值为 $1\Ums$、$5\Ums$、$9\Ums$、$\ldots$；若波向左传播，$v'=\frac{(n+\frac{3}{4})\lambda}{t_{2}-t_{1}}=(4n+3)\Ums$（$n=0,1,2,\ldots$），$v'$ 的可能值为 $3\Ums$、$7\Ums$、$11\Ums$、$\ldots$，故 ABC 正确，D 错误。
}

\item
%% number: 17
%% paperName: 2004年全国理综Ⅲ第 17 题 6 分
%% typeId: 1
%% type: 单选题
%% chapter: 曲线运动
%% point: 双星问题
%% method: 无
%% score: 6
%% degree: 450
%% duplicateId: 0
%% body:
我们的银河系的恒星中大约四分之一是双星。某双星由质量不等的星体 S$_{1}$ 和 S$_{2}$ 构成，两星在相互之间的万有引力作用下绕两者连线上某一定点 C 做匀速圆周运动。由天文观察测得其运动周期为 $T$，S$_{1}$ 到 C 点的距离为 $r_{1}$，S$_{1}$ 和 S$_{2}$ 的距离为 $r$，已知引力常量为 $G$。由此可求出 S$_{2}$ 的质量为\xzanswer{D}
\fourchoices[answer=D]
{$\frac{4\pi^{2}r^{2}(r-r_{1})}{GT^{2}}$}
{$\frac{4\pi^{2}r_{1}^{2}}{GT^{2}}$}
{$\frac{4\pi^{2}r^{2}}{GT^{2}}$}
{$\frac{4\pi^{2}r^{2}r_{1}}{GT^{2}}$}
%% answer: D
%% memo:
\memoanswer{
取 S$_{1}$ 为研究对象，由 $G\frac{m_{1}m_{2}}{r^{2}}=m_{1}\frac{4\pi^{2}}{T^{2}}r_{1}$，解得 $m_{2}=\frac{4\pi^{2}r^{2}r_{1}}{GT^{2}}$，故 D 正确。
}

\item
%% number: 18
%% paperName: 2004年全国理综Ⅲ第 18 题 6 分
%% typeId: 1
%% type: 单选题
%% chapter: 气体性质
%% point: 分子力
%% method: 无
%% score: 6
%% degree: 650
%% duplicateId: 0
%% body:
分子间有相互作用势能，规定两分子相距无穷远时两分子间的势能为零。设分子 a 固定不动，分子 b 以某一初速度从无穷远处向 a 运动，直到它们之间的距离最小。在此过程中，a、b 之间的势能\xzanswer{B}
\fourchoices[answer=B]
{先减小，后增大，最后小于零}
{先减小，后增大，最后大于零}
{先增大，后减小，最后小于零}
{先增大，后减小，最后大于零}
%% answer: B
%% memo:
\memoanswer{
当 b 分子以某一初速度从无穷远向 a 运动，在 $r>r_{0}$ 时，分子间表现为引力，分子力做正功，分子势能减小；当 $r<r_{0}$ 时，分子间表现为斥力，分子力做负功，分子势能增加，由于开始时 $E_{p}+E_{k}>0$，故分子动能为零时（即 a、b 相距最近时），分子势能大于零，故 B 正确。
}

\item
%% number: 19
%% paperName: 2004年全国理综Ⅲ第 19 题 6 分
%% typeId: 1
%% type: 单选题
%% chapter: 电磁感应
%% point: 电磁感应中图象
%% method: 无
%% score: 6
%% degree: 650
%% duplicateId: 0
%% body:
一矩形线圈位于一随时间 $t$ 变化的匀强磁场内，磁场方向垂直线圈所在的平面（纸面）向里，如图\ref{2004全国理综Ⅲ19a}所示。磁感应强度 $B$ 随 $t$ 的变化规律如图\ref{2004全国理综Ⅲ19b}所示。以 $I$ 表示线圈中的感应电流，以图\ref{2004全国理综Ⅲ19a}中线圈上箭头所示方向的电流为正，则以下的 $I-t$ 图中正确的是\xzanswer{A}
\twopicture[h=3.20cm, numA=1, numB=2, labelA=2004全国理综Ⅲ19a, labelB=2004全国理综Ⅲ19b]
{figs/19a.png}
{figs/19b.png}
\fourchoices[answer=A, ispicture=true, h=2.40cm]
{figs/19c.png}
{figs/19d.png}
{figs/19e.png}
{figs/19f.png}
%% answer: A
%% memo:
\memoanswer{
由题图可知，在 0~1s 时间内，磁感应强度均匀增大，感应电动势 $E=\frac{\Delta\Phi}{\Delta t}=S\cdot\frac{\Delta B}{\Delta t}$，感应电流 $I=\frac{E}{R}=\frac{S}{R}\cdot\frac{\Delta B}{\Delta t}$ 为定值，由楞次定律判断出感应电流方向为逆时针方向，与图\ref{2004全国理综Ⅲ19a}中电流相反，为负值；同理可知 3~4s 内与 0~1s 相同，1~2s 和 5~6s 内与 0~1s 电流等大反向，2~3s 和 4~5s 内磁通量不变，无感应电流，故 A 正确。
}

\item
%% number: 20
%% paperName: 2004年全国理综Ⅲ第 20 题 6 分
%% typeId: 1
%% type: 单选题
%% chapter: 光学
%% point: 透镜
%% method: 无
%% score: 6
%% degree: 450
%% duplicateId: 0
%% body:
如图所示，S 为一在 $xy$ 平面内的点光源。一平面镜垂直于 $xy$ 平面放置，它与 $xy$ 平面的交线为 MN，MN 与 $x$ 轴的夹角 $\theta=30^{\circ}$。现保持 S 不动，令平面镜以速率 $v$ 沿 $x$ 轴正方向运动，则 S 经平面镜所成的像\xzanswer{D}
\onepicture[h=3.80cm]{figs/20.png}
\fourchoices[answer=D]
{以速率 $v$ 沿 $x$ 轴正方向运动}
{以速率 $v$ 沿 $y$ 轴正方向运动}
{以速率 $\frac{1}{2}v$ 沿像与 S 连线方向向 S 运动}
{以速率 $v$ 沿像与 S 连线方向向 S 运动}
%% answer: D
%% memo:
\memoanswer{
由平面镜成像的规律和对称性可知，像沿 $S'S$ 连线向 S 运动，由几何关系得像的速度 $v_{0}=2v\sin30^{\circ}=v$，故 D 正确。
}

\item
%% number: 21
%% paperName: 2004年全国理综Ⅲ第 21 题 6 分
%% typeId: 1
%% type: 单选题
%% chapter: 电场
%% point: 带电粒子的圆周运动
%% method: 无
%% score: 6
%% degree: 250
%% duplicateId: 0
%% body:
一带正电的小球，系于长为 $l$ 的不可伸长的轻线一端，线的另一端固定在 O 点，它们处在匀强电场中，电场的方向水平向右，场强的大小为 $E$。已知电场对小球的作用力的大小等于小球的重力。现先把小球拉到图中的 P$_{1}$ 处，使轻线拉直，并与场强方向平行，然后由静止释放小球。已知小球在经过最低点的瞬间，因受线的拉力作用，其速度的竖直分量突变为零，水平分量没有变化，则小球到达与 P$_{1}$ 点等高的 P$_{2}$ 点时速度的大小为\xzanswer{B}
\onepicture[h=3.60cm]{figs/21.png}
\fourchoices[answer=B]
{$\sqrt{gl}$}
{$\sqrt{2gl}$}
{$2\sqrt{gl}$}
{$0$}
%% answer: B
%% memo:
\memoanswer{
小球到最低点时，水平速度满足 $v_{1}^{2}=2\frac{qE}{m}\cdot l$，即 $v_{1}=\sqrt{\frac{2qEl}{m}}$，从最低点到 P$_{2}$ 过程中，由动能定理得 $qEl-mgl=\frac{1}{2}mv_{2}^{2}-\frac{1}{2}mv_{1}^{2}$，其中 $qE=mg$，解得 $v_{2}=v_{1}=\sqrt{\frac{2Eql}{m}}=\sqrt{2gl}$，故 B 正确。
}

\item
%% number: 22
%% paperName: 2004年全国理综Ⅲ第 22 题 6 分
%% typeId: 4
%% type: 实验题
%% chapter: 直线运动
%% point: 游标卡尺和螺旋测微仪
%% method: 无
%% score: 6
%% degree: 850
%% duplicateId: 0
%% body:
用游标为 50 分度的卡尺（测量值可准确到 $0.02\Umm$）测定某圆筒的内径时，卡尺上的示数如图，可读出圆筒的内径为 \tkanswer[3.0cm]{54.14}$\Umm$。
\onepicture[align=right, w=6.50cm]{figs/22a.jpg}
%% answer: 
%% memo:
\memoanswer{
主尺读数为 $54\Umm$，游标尺读数为 $7\times0.02\Umm=0.14\Umm$，故读数为 $54\Umm+0.14\Umm=54.14\Umm$。
}

\item
%% number: 22
%% paperName: 2004年全国理综Ⅲ第 22 题 14 分
%% typeId: 4
%% type: 实验题
%% chapter: 电路
%% point: 测电动势与内阻
%% method: 无
%% score: 14
%% degree: 550
%% duplicateId: 0
%% body:
图中 $E$ 为直流电源，$R$ 为已知电阻，V 为理想电压表，其量程略大于电源电动势，K$_{1}$ 和 K$_{2}$ 为开关。现要利用图中电路测量电源的电动势 $E$ 和内阻 $r$，试写出主要实验步骤及结果表达式。
\onepicture[align=right, w=5.00cm]{figs/22b.jpg}
%% answer: 
\jdanswer{
\begin{enumerate}
	\item
	将 K$_{1}$ 闭合，K$_{2}$ 断开，记下电压表读数 $U_{1}$；
	\item
	将 K$_{1}$、K$_{2}$ 均闭合，记下电压表读数 $U_{2}$；
	\item
	结果表达式：电源电动势 $E=U_{1}$，内阻 $r=\frac{U_{1}-U_{2}}{U_{2}}R$。
\end{enumerate}
}
%% memo:
\memoanswer{
本题原卷及材料中未提供详解，待补充。
}

\item
%% number: 23
%% paperName: 2004年全国理综Ⅲ第 23 题 14 分
%% typeId: 6
%% type: 计算题
%% chapter: 运动和力
%% point: 连接体
%% method: 无
%% score: 14
%% degree: 650
%% duplicateId: 0
%% body:
如图所示，两个用轻线相连的位于光滑水平面上的物块，质量分别为 $m_{1}$ 和 $m_{2}$，拉力 $F_{1}$ 和 $F_{2}$ 方向相反，与轻线沿同一水平直线，且 $F_{1}>F_{2}$。试求在两个物块运动过程中轻线的拉力 $T$。
\onepicture[align=right, w=7.00cm]{\includesvg{figs/23}}
%% answer: 
\jdanswer{
$T=\frac{m_{1}F_{2}+m_{2}F_{1}}{m_{1}+m_{2}}$
}
%% memo:
\memoanswer{
设两物块一起运动的加速度为 $a$，由牛顿第二定律得

对整体，有

$F_{1}-F_{2}=(m_{1}+m_{2})a$ ①

对质量为 $m_{1}$ 的物块，有

$F_{1}-T=m_{1}a$ ②

由①②两式得 $T=\frac{m_{1}F_{2}+m_{2}F_{1}}{m_{1}+m_{2}}$ ③。
}

\item
%% number: 24
%% paperName: 2004年全国理综Ⅲ第 24 题 22 分
%% typeId: 6
%% type: 计算题
%% chapter: 磁场
%% point: 带电粒子在磁场中的运动
%% method: 无
%% score: 22
%% degree: 450
%% duplicateId: 0
%% body:
空间中存在方向垂直于纸面向里的匀强磁场，磁感应强度为 $B$，一带电量为 $+q$、质量为 $m$ 的粒子，在 P 点以某一初速开始运动，初速方向在图中纸面内如图中 P 点箭头所示。该粒子运动到图中 Q 点时速度方向与 P 点时速度方向垂直，如图中 Q 点箭头所示。已知 P、Q 间的距离为 $l$。若保持粒子在 P 点时的速度不变，而将匀强磁场换成匀强电场，电场方向与纸面平行且与粒子在 P 点时速度方向垂直，在此电场作用下粒子也由 P 点运动到 Q 点。不计重力。求：
\begin{enumerate}
	\item
	电场强度的大小；
	\item
	两种情况中粒子由 P 运动到 Q 点所经历的时间之差。
\end{enumerate}
\onepicture[align=right, w=4.00cm]{\includesvg{figs/24}}
%% answer: 
\jdanswer{
\begin{enumerate}
	\item
	$E=\sqrt{2}\frac{qB^{2}l}{m}$
	\item
	$t_{B}-t_{E}=\left(\frac{\pi}{2}-1\right)\frac{m}{qB}$
\end{enumerate}
}
%% memo:
\memoanswer{
（1）粒子在磁场中做匀速圆周运动，以 $v_{0}$ 表示粒子在 P 点的初速度，$R$ 表示圆周的半径，则有

$qv_{0}B=m\frac{v_{0}^{2}}{R}$ ①

由于粒子在 Q 点的速度垂直于它在 P 点时的速度，可知粒子由 P 点到 Q 点的轨迹为 $\frac{1}{4}$ 圆周，故有

$R=\frac{l}{\sqrt{2}}$ ②

以 $E$ 表示电场强度的大小，$a$ 表示粒子在电场中加速度的大小，$t_{E}$ 表示粒子在电场中由 P 点运动到 Q 点经过的时间，则有

$qE=ma$ ③

$R=\frac{1}{2}at_{E}^{2}$ ④

$R=v_{0}t_{E}$ ⑤

由以上各式，得

$E=\sqrt{2}\frac{qB^{2}l}{m}$ ⑥

（2）因粒子在磁场中由 P 点运动到 Q 点的轨迹为 $\frac{1}{4}$ 圆周，故运动经历的时间 $t_{B}$ 为圆周运动周期 $T$ 的 $\frac{1}{4}$，即有

$t_{B}=\frac{1}{4}T$ ⑦

而

$T=\frac{2\pi R}{v_{0}}$ ⑧

由⑦⑧和①式得

$t_{B}=\frac{\pi m}{2qB}$

由①⑤两式得 $t_{E}=\frac{m}{qB}$ ⑨

故经历的时间差为 $t_{B}-t_{E}=\left(\frac{\pi}{2}-1\right)\frac{m}{qB}$。
}

\item
%% number: 25
%% paperName: 2004年全国理综Ⅲ第 25 题 22 分
%% typeId: 6
%% type: 计算题
%% chapter: 运动和力
%% point: 动量和能量
%% method: 无
%% score: 22
%% degree: 350
%% duplicateId: 0
%% body:
如图所示，在一光滑的水平面上有两块相同的木板 B 和 C。重物 A（视为质点）位于 B 的右端，A、B、C 的质量相等。现 A 和 B 以同一速度滑向静止的 C、B 与 C 发生正碰。碰后 B 和 C 粘在一起运动，A 在 C 上滑行，A 与 C 有摩擦力。已知 A 滑到 C 的右端而未掉下。试问：从 B、C 发生正碰到 A 刚移到 C 右端期间，C 所走过的距离是 C 板长度的多少倍。
\onepicture[align=right, w=9.00cm]{\includesvg{figs/25}}
%% answer: 
\jdanswer{
$\frac{7}{3}$
}
%% memo:
\memoanswer{
设 A、B、C 的质量均为 $m$，碰撞前，A 与 B 的共同速度为 $v_{0}$，碰撞后 B 与 C 的共同速度为 $v_{1}$。对 B、C，由动量守恒定律得

$mv_{0}=2mv_{1}$ ①

设 A 滑至 C 的右端时，三者的共同速度为 $v_{2}$。对 A、B、C，由动量守恒定律得

$2mv_{0}=3mv_{2}$ ②

设 A 与 C 的动摩擦因数为 $\mu$，从发生碰撞到 A 移至 C 的右端时 C 所走过的距离为 $s$，对 B、C，由功能关系

$\mu mgs=\frac{1}{2}(2m)v_{2}^{2}-\frac{1}{2}(2m)v_{1}^{2}$ ③

设 C 的长度为 $l$，对 A，由功能关系

$\mu mg(s+l)=\frac{1}{2}mv_{0}^{2}-\frac{1}{2}mv_{2}^{2}$ ④

由以上各式解得

$\frac{s}{l}=\frac{7}{3}$ ⑤
}

\end{enumerate}

\end{document}
